\begin{afterword}
\summaryhead

The post opens with a joke. A physics instructor shows students a metal plate beside a radiator. The side near the radiator feels cool and the far side warm. The students propose convection, strange metals and other causes, and none says ``I don't know.'' The instructor had turned the plate around.

The author asks whether an answer like ``Eh, maybe because of the heat conduction and so?'' is a real belief. Heat conduction normally predicts a warmer near side; if the phrase can also explain a cooler one, it can explain anything, and so it is no better than ``magic.'' The two phrases differ only in literary genre: one sounds like \textsc{Star Trek}, the other like \textsc{Buffy the Vampire Slayer}. A physicist would instead measure temperatures, use the diffusion equation, and conclude that the plate had been turned around. The students' deeper error, the post concludes, was to think that science-sounding words are physics.

\responsehead

The post makes a sound point with a good joke: a phrase that could explain any outcome explains none, and sounding scientific does not change that. Two things limit it.

First, the point is not new, and the post does not say so. Richard Feynman made it with the same device in 1985. A textbook that answered ``Energy makes it go'' for every picture, he wrote, might as well have said ``Wakalixes makes it go,'' because ``There's no knowledge coming in.'' The principle behind both, that an account which fits any outcome says nothing, is Karl Popper's. What the post adds is the physics-classroom example and the idea of a literary genre. Its restatement also adds two technical terms, ``anticipation-isomorphic'' and ``a disguised hypothesis of maximum entropy,'' for the plain claim that the answer rules nothing out; neither is defined.

Second, the central sentence says more than the example shows. The post claims that if ``because of heat conduction'' can explain the near side feeling cooler, it can explain ``pretty much anything.'' But heat conduction, with the right starting condition, does explain a cooler near side. That is the physicist's answer in the same post: a plate turned around and returning to equilibrium. The student's answer fails for another reason: the student names no starting condition that would produce the observation and calculates nothing. The qualifier ``used in such fashion'' points this way, but the comparison with ``magic'' treats the two phrases as equals. They are not: one comes with an equation, which the post itself brings in when it describes measuring the plate, and the other does not.

This affects what a reader takes away. The fault the example shows is saying the phrase instead of using the equation. The next post, ``Guessing the Teacher's Password,'' grants as much, with a ``Maybe'': ``heat conduction?'' can be a first step if it leads to a calculation.

In short: Feynman's ``wakalixes'' point, uncredited, applied to a joke whose own answer shows that heat conduction explains the observation once the missing starting condition is supplied.
\end{afterword}
